In the kinetic description, electrons are described by a distribution
function $f\left( {\bf r},{\bf p},t\right) $, which gives the density in
phase space of particles with a momentum ${\bf r}$ at a position ${\bf r}$
and at time $t$. The particle conservation equation in phase space is the
Boltzmann equation 
\begin{equation}
\frac{\partial f}{\partial t}+{\bf v}\cdot \nabla _{{\bf r}}f+q_{e}\left[ 
{\bf E}\left( {\bf r},t\right) +{\bf v}\times {\bf B}\left( {\bf r},t\right)
\right] \cdot \nabla _{{\bf p}}f=\left. \frac{\partial f}{\partial t}\right|
_{\text{C}}  \label{kin1}
\end{equation}
where 
\begin{equation}
\left. \frac{\partial f}{\partial t}\right| _{\text{C}}\equiv C\left(
f\right)
\end{equation}
is the collision operator. The fields ${\bf E}\left( {\bf r},t\right) $ and $%
{\bf B}\left( {\bf r},t\right) $ are assumed to consist of time-independent
macroscopic fields ${\bf E}_{0}\left( {\bf r}\right) $ and ${\bf B}%
_{0}\left( {\bf r}\right) $ and fields associated with plane waves. 
\begin{eqnarray}
{\bf E}\left( {\bf r},t\right) &=&{\bf E}_{0}\left( {\bf r}\right) +\sum_{k}%
{\bf E}_{k}e^{i\left( {\bf k}\cdot {\bf r}-\omega t\right) } \\
{\bf B}\left( {\bf r},t\right) &=&{\bf B}_{0}\left( {\bf r}\right) +\sum_{k}%
{\bf B}_{k}e^{i\left( {\bf k}\cdot {\bf r}-\omega t\right) }
\end{eqnarray}

Because we are interested in solving the kinetic equation on the bounce and
collisional time scales, we need to average over the faster time scales,
which are the gyromotion and the wave oscillation.

Performing a time-averaging $\int_{0}^{2\pi /\omega }dt$ of the equation (%
\ref{kin1}) removes the fast wave time scale from the equation, to give 
\begin{equation}
\frac{\partial \widetilde{f}}{\partial t}+{\bf v}\cdot \nabla _{{\bf r}}%
\widetilde{f}+q_{e}\left[ {\bf E}_{0}\left( {\bf r}\right) +{\bf v}\times 
{\bf B}_{0}\left( {\bf r}\right) \right] \cdot \nabla _{{\bf p}}\widetilde{f}%
=\left. \frac{\partial \widetilde{f}}{\partial t}\right| _{\text{C}%
}-\int_{0}^{2\pi /\omega }dt\;\sum_{k}\left( q_{e}\left[ {\bf E}_{k}+{\bf v}%
\times {\bf B}_{k}\right] \cdot \nabla _{{\bf p}}f\right)  \label{kin2}
\end{equation}
where 
\begin{equation}
\widetilde{f}=\int_{0}^{2\pi /\omega }dt\;f
\end{equation}
is the wave-period averaged distribution function. The time derivative in
the first term of (\ref{kin2}) implicitely refers to times longer than the
wave period.

Under the assumption of a strong magnetic field, such that the gyrofrequency 
$\Omega _{e}$%
\begin{equation}
\Omega _{e}=\frac{q_{e}B_{0}}{\gamma m_{e}}
\end{equation}
is much larger than both the collisional frequency and the bounce frequency,
we can expand the distribution function 
\begin{equation}
\widetilde{f}=\widetilde{f}_{0}+\widetilde{f}_{1}+\widetilde{f}_{2}+\cdots
\end{equation}
with a small parameter 
\begin{equation}
\delta \sim \frac{\omega _{b}}{\Omega _{e}}\sim \frac{\nu _{c}}{\Omega _{e}}
\end{equation}

The zero order equation becomes 
\begin{equation}
q_{e}{\bf v}\times {\bf B}_{0}\left( {\bf r}\right) \cdot \nabla _{{\bf p}}%
\widetilde{f}_{0}=0  \label{kinexp0}
\end{equation}

We have, in the $\left( p_{\Vert },p_{\bot },\varphi \right) $ space 
\begin{equation}
{\bf v}=\frac{{\bf p}}{\gamma \left( p_{\Vert },p_{\bot }\right) m}
\end{equation}
with the momentum being (\ref{pparperp}) 
\begin{equation}
{\bf p}=p_{\Vert }\widehat{\Vert }+p_{\bot }\widehat{\bot }
\end{equation}
and the gradient (\ref{gradparperp}) 
\begin{equation}
\nabla _{{\bf p}}\widetilde{f}=\frac{\partial \widetilde{f}}{\partial
p_{\Vert }}\widehat{\Vert }+\frac{\partial \widetilde{f}}{\partial p_{\bot }}%
\widehat{\bot }+\frac{1}{p_{\bot }}\frac{\partial \widetilde{f}}{\partial
\varphi }\widehat{\varphi }  \label{gradparperp2}
\end{equation}

In this system, by definition, 
\begin{equation}
{\bf B}_{0}\left( {\bf r}\right) =B_{0}\left( {\bf r}\right) \widehat{\Vert }
\end{equation}
so that the gyromotion operator becomes 
\begin{eqnarray}
q_{e}{\bf v}\times {\bf B}_{0}\left( {\bf r}\right) \cdot \nabla _{{\bf p}}
&=&\Omega _{e}\left( {\bf p}\times \widehat{\Vert }\right) \cdot \nabla _{%
{\bf p}}\widetilde{f}  \nonumber \\
&=&-\Omega _{e}\frac{\partial }{\partial \varphi }  \label{gyromotion}
\end{eqnarray}

The equation (\ref{kinexp0}) becomes 
\begin{equation}
\frac{\partial \widetilde{f}_{0}}{\partial \varphi }=0  \label{ftilde0phi}
\end{equation}
and therefore $\widetilde{f}_{0}$ is independent of $\varphi $. The first
order equation is 
\begin{equation}
\frac{\partial \widetilde{f}_{0}}{\partial t}+{\bf v}\cdot \nabla _{{\bf r}}%
\widetilde{f}_{0}+q_{e}{\bf E}_{0}\left( {\bf r}\right) \cdot \nabla _{{\bf p%
}}\widetilde{f}_{0}+q_{e}{\bf v}\times {\bf B}_{0}\left( {\bf r}\right)
\cdot \nabla _{{\bf p}}\widetilde{f}_{1}=\left. \frac{\partial \widetilde{f}%
_{0}}{\partial t}\right| _{\text{C}}-\int_{0}^{2\pi /\omega
}dt\;\sum_{k}\left( q_{e}\left[ {\bf E}_{k}+{\bf v}\times {\bf B}_{k}\right]
\cdot \nabla _{{\bf p}}f\right)  \label{kinexp}
\end{equation}

The last term in the equation (\ref{kinexp}) has been calculated by Lerche
for a uniform plasma, in the form of a quasilinear operator ${\cal Q}\left( 
\overline{f}\right) $. We rewrite 
\begin{eqnarray}
\;\left. \frac{\partial \widetilde{f}_{0}}{\partial t}\right| _{\text{C}} &=&%
{\cal C}\left( \widetilde{f}_{0}\right) \\
-\;\int_{0}^{2\pi /\omega }dt\;\sum_{k}\left( q_{e}\left[ {\bf E}_{k}+{\bf v}%
\times {\bf B}_{k}\right] \cdot \nabla _{{\bf p}}f\right) &=&{\cal Q}\left( 
\widetilde{f}_{0}\right)
\end{eqnarray}

Performing the gyro-averaging $\int_{0}^{2\pi }d\varphi $ on the kinetic
equation (\ref{kinexp}), we find, using (\ref{ftilde0phi}), that 
\begin{equation}
\int_{0}^{2\pi }d\varphi \;\frac{\partial \widetilde{f}_{0}}{\partial t}=%
\frac{\partial \widetilde{f}_{0}}{\partial t}
\end{equation}
and 
\begin{equation}
\int_{0}^{2\pi }d\varphi \;{\bf v}\cdot \nabla _{{\bf r}}\widetilde{f}%
_{0}=\int_{0}^{2\pi }\left( d\varphi \;{\bf v}\right) \cdot \nabla _{{\bf r}}%
\widetilde{f}_{0}{\bf =v}_{\text{gc}}\cdot \nabla _{{\bf r}}\widetilde{f}_{0}
\end{equation}

Concerning the electric field, we decompose the gradient in momentum space
using (\ref{gradparperp2}) 
\begin{eqnarray}
\int_{0}^{2\pi }d\varphi \;q_{e}{\bf E}_{0}\left( {\bf r}\right) \cdot
\nabla _{{\bf p}}\widetilde{f}_{0} &=&q_{e}{\bf E}_{0}\left( {\bf r}\right)
\cdot \int_{0}^{2\pi }d\varphi \left[ \frac{\partial \widetilde{f}_{0}}{%
\partial p_{\Vert }}\widehat{\Vert }+\frac{\partial \widetilde{f}_{0}}{%
\partial p_{\bot }}\widehat{\bot }+\frac{1}{p_{\bot }}\frac{\partial 
\widetilde{f}_{0}}{\partial \varphi }\widehat{\varphi }\right] \\
&=&q_{e}\frac{\partial \widetilde{f}_{0}}{\partial p_{\Vert }}{\bf E}%
_{0}\left( {\bf r}\right) \cdot \widehat{\Vert } \\
&&+q_{e}\frac{\partial \widetilde{f}_{0}}{\partial p_{\bot }}{\bf E}%
_{0}\left( {\bf r}\right) \cdot \int_{0}^{2\pi }d\varphi \widehat{\bot } \\
&&+\frac{q_{e}}{p_{\bot }}{\bf E}_{0}\left( {\bf r}\right) \cdot
\int_{0}^{2\pi }d\varphi \frac{\partial \widetilde{f}_{0}}{\partial \varphi }%
\widehat{\varphi }
\end{eqnarray}
and, using 
\begin{eqnarray}
\int_{0}^{2\pi }d\varphi \widehat{\bot } &=&0 \\
\int_{0}^{2\pi }d\varphi \frac{\partial \widetilde{f}_{0}}{\partial \varphi }%
\widehat{\varphi } &=&-\widetilde{f}_{0}\int_{0}^{2\pi }d\varphi \frac{%
\partial \widehat{\varphi }}{\partial \varphi }=\widetilde{f}%
_{0}\int_{0}^{2\pi }d\varphi \widehat{\bot }=0
\end{eqnarray}
we obtain 
\[
\int_{0}^{2\pi }d\varphi \;q_{e}{\bf E}_{0}\left( {\bf r}\right) \cdot
\nabla _{{\bf p}}\widetilde{f}_{0}=q_{e}E_{0\Vert }\left( {\bf r}\right) 
\frac{\partial \widetilde{f}_{0}}{\partial p_{\Vert }} 
\]

The gyromotion term is averaged to zero 
\begin{equation}
\int_{0}^{2\pi }d\varphi \;q_{e}{\bf v}\times {\bf B}_{0}\left( {\bf r}%
\right) \cdot \nabla _{{\bf p}}\widetilde{f}_{1}=-\Omega _{e}\int_{0}^{2\pi %
}d\varphi \;\frac{\partial \widetilde{f}_{1}}{\partial \varphi }=0
\end{equation}
so that we get finally 
\begin{equation}
\frac{\partial \widetilde{f}_{0}}{\partial t}+{\bf v}_{\text{gc}}\cdot %
\nabla _{{\bf r}}\widetilde{f}_{0}+\frac{\partial }{\partial p_{\Vert }}%
q_{e}E_{0\Vert }\left( {\bf r}\right) \widetilde{f}_{0}={\cal C}\left( 
\widetilde{f}_{0}\right) +{\cal Q}\left( \widetilde{f}_{0}\right) 
\end{equation}

This equation is called drift-kinetic equation. Renaming the guiding-center
distribution function $\widetilde{f}_{0}=f\left( {\bf r},p_{\Vert },p_{\bot
},t\right) $ we get 
\begin{equation}
\frac{\partial f}{\partial t}+{\bf v}_{\text{gc}}\cdot \nabla f={\cal C}%
\left( f\right) +{\cal Q}\left( f\right) +{\cal E}\left( f\right) 
\end{equation}
where we define an electric field operator
\begin{equation}
{\cal E}\left( f\right) =-\frac{\partial }{\partial p_{\Vert }}%
q_{e}E_{0\Vert }\left( {\bf r}\right) f
\end{equation}
